\begin{corollary}[Instantaneous kernel-loss gradient space]
\label{cor:gradientspace}
Under Theorem~\ref{thm:gram}, set $g(\theta)=\operatorname{vech}G(\theta)$
and $J_g=D_\theta g$. For a differentiable loss depending on $\theta$
only through finitely many fixed-input kernel evaluations,
\begin{equation}
 \nabla_\theta L\in\operatorname{im}J_g(\theta)^T
       =(\ker J_g(\theta))^\perp.
\end{equation}
This space has dimension $\operatorname{rank}J_g(\theta)\le r$,
with equality on the regular parameter set.
\end{corollary}
\begin{proof}
Each evaluated kernel is a smooth function of $g$, so the loss factors
as $L=\ell\circ g$. The chain rule gives
$\nabla_\theta L=J_g^T\nabla_g\ell$; the orthogonal-complement
identity follows from finite-dimensional linear algebra.
\end{proof}
This is an instantaneous constraint, not confinement to one fixed
$r$-dimensional linear subspace along training. The subspace can rotate
with $\theta$; finite steps, momentum and preconditioning need not stay
in it. Direct parameter regularization can add components outside it.
Nor does this imply $r$ circuit evaluations suffice for a gradient.
